\section{D. Points, central extensions, and arithmetic realizations}
\label{sec:D}
Here $M$ means exactly the scheme of C3-399(13).
All extension sequences of group schemes are fppf exact.
This is a mathematical reconstruction; it is not a claim
that the inaccessible original files were recovered byte for byte.

\begin{statement}{D.1}{Proved}{The exact all-ring functor}
For every commutative ring $R$, put

$$
x(a,b)=\begin{pmatrix}a&b\\-b&a+b\end{pmatrix},\quad
 n=a^2+ab+b^2,\quad
 y(c,d,e)=\begin{pmatrix}c&d\\-ed&ec\end{pmatrix},\quad q=c^2+d^2.
$$

Then $M(R)$ is exactly the triples $(x,y,f)$ satisfying

$$
e^2=t^2=1,\quad q\in R^\times,\quad n=etq,\quad
 f\in\operatorname{GL}_2(R),\quad\det f=eq.
$$

If $H$ is the functor of the six scalar coordinates and
$\delta=eq$, then

$$
M\simeq\operatorname{SL}_2\rtimes H,\qquad
 (g,h)\longmapsto(x(h),y(h),g\operatorname{diag}(\delta(h),1)).
$$

The action is conjugation by $\operatorname{diag}(\delta,1)$.
\end{statement}
\begin{proof}
The determinants are $n$ and $eq$.
The source relation $\det f=\det y=t\det x$ is equivalent
to the equations above, since $t^2=1$.
Multiplication of scalar coordinates is

$$
(a,b)(a',b')=(aa'-bb',ab'+ba'+bb'),
$$


$$
(c,d,e)(c',d',e')=(cc'-e'dd',cd'+e'dc',ee'),\quad t''=tt'.
$$

These polynomial laws make $H$ a group and $\delta$ a character.
Factoring $f$ uniquely as $g\operatorname{diag}(\delta,1)$
proves the semidirect product, including over nilpotent rings.
This is also the point description of B.1.
\end{proof}

\begin{statement}{D.2}{Proved}{Integral, local, and adelic points}
The scalar coordinates are precisely the norm fibres

$$
x\in N_\rho^{-1}(etq),\qquad c-di\in N_\sigma^{-1}(q)
$$

in $R[T]/(T^2-T+1)$ and $R[i]$.
For $R=\mathbb Z_p$ the necessary and sufficient restrictions are

$$
p=2:\ q\equiv1\pmod4;\qquad
 p=3:\ q\equiv et\pmod3;\qquad p\ge5:\ \text{none}.
$$

For $\mathbb Q_p$, write $q=p^ku$ with $u$ a unit. Then

$$
\begin{array}{c|l}
p&\text{conditions}\\\hline
2&k\text{ even},\ u\equiv1\pmod4\\
3&k\text{ even},\ u\equiv et\pmod3\\
p\ge5,\ p\equiv1\pmod{12}&\text{none}\\
p\ge5,\ p\not\equiv1\pmod{12}&k\text{ even}.
\end{array}
$$

Moreover

$$
M(\mathbb Z)=(C_6\times C_4)\rtimes\operatorname{GL}_2(\mathbb Z),
 \quad M(\widehat{\mathbb Z})=\prod_pM(\mathbb Z_p),\quad
 M(\mathbb A_f)=\prod_p'M(\mathbb Q_p).
$$

Determinant minus one acts by inversion on $C_4$ and trivially
on $C_6$. For $M^+=\ker\delta$, impose $q=e$;
this forces $e=1$ at two and $t=1$ at three.
In particular $M^+(\mathbb Z)=C_6\times C_4\times\operatorname{SL}_2(\mathbb Z)$.
\end{statement}
\begin{proof}
Split quadratic algebras have all scalars as norms.
An unramified quadratic extension has all unit norms and even
norm valuations: finite-field norm surjectivity lifts through
successive powers of the maximal ideal using the surjective trace.
The Gaussian algebra ramifies only at two.
Its unit norms are one modulo four; conversely a unit one
modulo eight is a square, and a unit five modulo eight is
$(\sqrt{u-4})^2+2^2$. Also $N(1+i)=2$.
The Eisenstein algebra ramifies only at three.
Its unit norms reduce to squares, hence to one; every such
unit is a square by Hensel. Also $N(2\rho-1)=3$.
Combining both norm conditions gives the tables.
Both quadratic algebras split at $p\ge5$ exactly for
$p\equiv1\pmod{12}$.

Over $\mathbb Z$, positivity forces both invertible norms
to equal one, giving unit groups $C_6,C_4$ and $e=t=\det f$.
The section $f\mapsto(1,\operatorname{diag}(1,\det f),f)$
has the indicated action. Affine functors preserve products
of rings. Finite presentation gives the restricted product
over the finite adeles, using these integral subgroups.
Substituting $q=e$ proves the determinant-one assertions.
\end{proof}

\begin{statement}{D.3}{Refuted}{The literal cyclotomic target map is impossible}
There is no continuous homomorphism
$G_{\mathbb Q}\to M(\widehat{\mathbb Z})$ with cyclotomic
determinant. A multiplier-preserving homomorphism from
$\widehat{GT}$ or $\widehat{GT}_0$ therefore cannot restrict
to this proposed arithmetic map.
This concerns this exact $M$, not a different normalizer
or an extension of $G_{\mathbb Q}$.
\end{statement}
\begin{proof}
The projections to $T_\rho(\mathbb Z_5)$ and to the sign $t$
are abelian. By Kronecker--Weber they factor through
$G_{\mathbb Q}^{ab}=\widehat{\mathbb Z}^{\times}$ via the
cyclotomic character \cite[I, Theorem 4.16]{D_MilneCFT}.
Choose the order-four element whose fifth component is the
Teichmuller lift of two and whose other components are one.
Its image in $T_\rho(\mathbb F_5)=\mathbb F_{25}^{\times}$
has order dividing four, so its norm, its sixth power,
is $1$ or $4$. The equation $N_\rho(x)=t\chi_5$ instead
requires $2$ or $3$. This is a contradiction.
The executed enumeration finds exactly the four fourth-torsion
pairs $(a,b)=(1,0),(2,0),(3,0),(4,0)$, with norms $1,4,4,1$.
Composition with the arithmetic inclusion proves the GT
consequence.
\end{proof}

\begin{statement}{D.4}{Refuted}{No arithmetic image is open}
Every continuous homomorphism
$G_{\mathbb Q}\to M(\widehat{\mathbb Z})$ has nonopen image.
No determinant hypothesis is needed.
\end{statement}
\begin{proof}
At seven, $T_\rho(\mathbb Z_7)=(\mathbb Z_7^\times)^2$.
The projection of $M(\widehat{\mathbb Z})$ to this group is
onto: choose $e=t=1$, a Gaussian unit of the prescribed norm,
and $f$ of that determinant, taking identity elsewhere.
It is an open quotient of compact groups.
The arithmetic projection factors through
$\prod_p\mathbb Z_p^\times$.
After applying the two seven-adic logarithms, every factor
$p\ne7$ vanishes: its open pro-$p$ subgroup and its finite
quotient have zero image in the torsion-free pro-seven group.
Continuity and density of finite-support elements kill their
product. The remaining factor has rank one, whereas the
target principal units have rank two.
Modulo $7^r$ the index is at least $7^{r-1}$, hence unbounded.
\end{proof}

\begin{statement}{D.5}{Proved}{Finite GT equation shadows}
Modulo $n$, put

$$
X=\begin{pmatrix}-1&2\\0&-1\end{pmatrix},\
 Y=\begin{pmatrix}-1&0\\-2&-1\end{pmatrix},\
 Z=(XY)^{-1},\ S=\sigma,\ R=Y^{-1}\rho.
$$

Let $H_n=\langle X,Y\rangle$ have exponent $d_n$.
For $\lambda\in(\mathbb Z/2d_n)^\times$, $m=(\lambda-1)/2$,
and $F\in[H_n,H_n]$, impose

$$
SFS^{-1}F=1,\qquad
 (R^2FR^{-2})Z^m(RFR^{-1})Y^mFX^m=1.
$$

The exhaustive counts are

$$
\begin{array}{c|rrrrr}
n&2&3&4&6&12\\\hline
|H_n|&1&24&4&24&96\\
|H_n'|&1&8&1&8&8\\
d_n&1&12&2&12&12\\
\text{pairs}&1&24&2&24&24.
\end{array}
$$

Every counted pair gives an automorphism
$X\mapsto X^\lambda$, $Y\mapsto F^{-1}Y^\lambda F$.
These are finite shadows, not asserted global images.
\end{statement}
\begin{proof}
Conjugation by $S$ exchanges $X,Y$, and conjugation by $R$
cycles $X,Y,Z$. Thus these are the finite two-cycle and
three-cycle relations of \cite[Section 0.1]{D_HS}.
The program enumerates $H_n$, closes all commutators,
computes every element order, and tests every allowed pair.
It verifies the proposed homomorphisms on the complete
Cayley graph with $X,Y,X^{-1},Y^{-1}$ edges, checking
consistency at every revisited vertex and distinctness
of all images. This proves bijective homomorphisms,
without assuming an unverified finite presentation.
Compatibility across quotients and the pentagon are
additional conditions.
\end{proof}

\begin{statement}{D.6}{Proved}{A noncongruence quotient}
The kernel

$$
K=\ker\bigl(\widehat{\operatorname{SL}_2(\mathbb Z)}
       \to\operatorname{SL}_2(\widehat{\mathbb Z})\bigr)
$$

surjects onto $\operatorname{SL}_2(\mathbb F_8)$.
This does not construct a GT automorphism.
\end{statement}
\begin{proof}
For $\alpha^3+\alpha+1=0$, set

$$
A=\begin{pmatrix}0&1\\1&0\end{pmatrix},\quad
 B=\begin{pmatrix}0&\alpha\\\alpha^{-1}&1\end{pmatrix}.
$$

Their orders are $2,3$, and $AB$ has order $7$.
The amalgam gives a homomorphism $\sigma\mapsto A,\rho\mapsto B$.
The program enumerates all $504$ elements. The conjugacy
class orders and sizes are

$$
(1,1),(2,63),(3,56),(7,72),(7,72),(7,72),
 (9,56),(9,56),(9,56).
$$

The normal closure of each nonidentity class is all $504$
elements, proving simplicity by exhaustive calculation.

The image of $K$ is normal, hence trivial or all.
If trivial, this simple quotient would factor through a
finite product of $\operatorname{SL}_2(\mathbb Z/p^a)$.
One factor would surject. Its congruence kernel is a
$p$-group and must have trivial normal image.
For $p=2,3,7$, the residual groups have orders $6,24,336$.
For any other $p$, the target has no element of order $p$;
its elementary unipotent subgroups would be killed, but
they generate $\operatorname{SL}_2(\mathbb F_p)$.
All cases contradict surjectivity.
\end{proof}

\begin{statement}{D.7}{Proved}{An extra adelic sign}
There is a determinant-one point with $e_5=-1$ and
$e_p=1$ elsewhere. It violates $e=\varepsilon_2(\delta)$.
\end{statement}
\begin{proof}
Choose $i_5^2=-1$ in $\mathbb Z_5$ and take at five
$x=I$, $y=\operatorname{diag}(i_5,-i_5)$, $f=I$,
$e=-1$, $t=1$. Then $q=-1$ and both required determinants
are one. Use identities elsewhere.
The checked residues of $i_5$ modulo $5^1,\ldots,5^6$
are $2,7,57,182,2057,14557$.
Every character takes value one at $\delta=1$,
whereas this $e$ does not.
\end{proof}

\begin{statement}{D.8}{Open}{The corrected global comparison remains unfinished}
After a corrected target or source extension is specified,
the requested global image/kernel comparison and an explicit
GT automorphism invisible to a congruence representation
remain unproved here. The missing finite-to-global step is
compatible lifting through all finite quotients and, for
$\widehat{GT}$, the pentagon. D.6 is not that step.
\end{statement}
\begin{proof}[Proved cases and limitation]
The pairs $(1,1)$ and $(-1,1)$ are global examples.
The second inverts both free profinite generators and is
an involution. Its two-cycle and pentagon are trivial,
and its three-cycle is $z^{-1}y^{-1}x^{-1}=1$.
With standard real tangential basepoints it is complex
conjugation. The definitions and arithmetic inclusion
are those in \cite[Section 0.1 and Main Theorem]{D_HS}.
D.5 proves no lifting theorem for the other finite pairs.
For the literal target the obstruction already occurs
modulo five in D.3.
\end{proof}

\begin{statement}{D.9}{Proved}{Marked central extensions of the Klein four group}
For $V=C_2^2$, marked central extensions by $C_2$
form $H^2(V,C_2)=\mathbb F_2^3$.
Write $X^2=z^a$, $Y^2=z^b$, $[X,Y]=z^c$.
For $c=0$ the group is $C_2^3$ when $a=b=0$,
and $C_4\times C_2$ otherwise.
For $c=1$ it is $Q_8$ when $a=b=1$, and $D_8$ otherwise.
The square function is $au+bv+cuv$.
\end{statement}
\begin{proof}
Changing lifts by $z$ changes none of these invariants.
Every element has normal form $z^wX^uY^v$.
Conversely the law

$$
(w,u,v)(w',u',v')=
 (w+w'+auu'+bvv'+cvu',u+u',v+v')
$$

is associative by bilinearity and realizes every triple.
It also shows addition of parameters under Baer sum.
Squaring gives the stated quadratic form.
For $c=1$ the centre is $\langle z\rangle$; all six
noncentral elements have order four precisely in the
quaternion case. Otherwise an involution and an order-four
element give the dihedral presentation.
The independent cochain computation tests $512$ normalized
cochains, obtaining $16$ cocycles, $2$ coboundaries,
and $8$ classes.
\end{proof}

\begin{statement}{D.10}{Proved}{The manuscript square rule}
Writing $\varepsilon_3=(-1)^u$, $\varepsilon_2=(-1)^v$,
the rule that the square is one exactly when
$\varepsilon_3=1$ or $\varepsilon_3\varepsilon_2=1$
gives $q(u,v)=u+uv$. Its class is $(1,0,1)$, hence $D_8$.
Both isotropic lines have split inverse-image extensions,
but their order-two complements are nonnormal in the
whole group. Pullback gives $\alpha^2+\alpha\beta$.
\end{statement}
\begin{proof}
Use the cocycle $u(u'+v')$.
Its commutator pairing $uv'+u'v$ is nondegenerate.
For any nonzero isotropic $\ell$, choose $w$ pairing to one.
Conjugation sends an involutive lift $L$ to $zL$,
which is outside $\{1,L\}$. Pullback gives the stated cup class.

This retains the logical OR on C3-602.
The separate reflection ambiguity on C3-608 is not used:
the exchange matrix conjugates $\rho$ to $\rho^{-1}$,
whereas $\operatorname{diag}(1,-1)$ does not.
The order and kernel ambiguities on scans 609--610
remain separately recorded in the readings register.
\end{proof}

\begin{statement}{D.11}{Refuted}{The literal class on C3-611 is zero}
For the quadratic characters literally displayed there,

$$
\chi_{-3}\smile\chi_3=(-3,3)=0
       \quad\hbox{in }\operatorname{Br}(\mathbb Q)[2].
$$

The characters belong to $H^1$; the cup belongs to $H^2$.
With $\alpha=\chi_{-3}$, $\beta=\chi_{-1}$ this is
D.10's $\alpha(\alpha+\beta)$.
\end{statement}
\begin{proof}
The matrices

$$
I=\begin{pmatrix}0&-3\\1&0\end{pmatrix},\qquad
 J=\begin{pmatrix}0&3\\1&0\end{pmatrix}
$$

satisfy $I^2=-3$, $J^2=3$, $IJ=-JI$.
The four matrices $1,I,J,IJ$ are independent and span
$M_2(\mathbb Q)$, proving splitting.
Equivalently $3$ is the norm of $\sqrt{-3}$.
The product of square classes $(-3)(-1)=3$ gives
the identification with D.10.
Thus the literal class has order one, not a nonzero
Brauer order. This is Task D.d and Lead 4.
\end{proof}

\begin{statement}{D.12}{Proved}{The distinct cyclic-four obstruction}
Pulling $1\to C_2\to C_4\to C_2\to1$ back by $\chi_{-3}$
instead gives

$$
\chi_{-3}^2=(-3,-3)=(-1,-3).
$$

Its Brauer order is two and its ramified places are
$3,\infty$. It splits over $\mathbb Q(i)$ and
$\mathbb Q(\sqrt{-3})$, but not over $\mathbb Q(\sqrt3)$.
\end{statement}
\begin{proof}
The cyclic-four cocycle is $uu'$, giving the cup square;
the quaternion identity $(a,a)=(a,-1)$ gives the symbol.
For a field of characteristic not two, $(-1,b)$ splits
exactly when $b$ is a norm from $F(i)$.
If $N(z)=b$, multiplication by $i$ and
$v\mapsto z\bar v$ give splitting matrices.
Conversely a split module is rank one over $F(i)$,
and its anticommuting operator has this form.

At infinity, $-3$ is not a norm.
At three the Gaussian extension is unramified, so its
odd valuation excludes a norm.
At two, $-3=(\sqrt{-7})^2+2^2$, with $-7\equiv1\pmod8$.
At every other prime the algebra is split or $-3$ is a
unit norm in an unramified quadratic extension.
Quaternion conjugation identifies the algebra with its
opposite, so the nonzero class has order two.
Making a parameter square splits it.
A real embedding of $\mathbb Q(\sqrt3)$ retains the real
obstruction. The $C_4$ notation in C3-631 describes the
extension group, not the Brauer order.
\end{proof}

\begin{statement}{D.13}{Refuted}{Imaginary quadratic quotients of constant Dic3}
For

$$
\operatorname{Dic}_3=C_3\rtimes C_4
 =\langle r,s:r^3=s^4=1,\ srs^{-1}=r^{-1}\rangle,
$$

no homomorphism from $G_{\mathbb Q}$ induces an imaginary
quadratic character under its index-two quotient.
This does not settle a cocycle for an unspecified
nonconstant group scheme.
\end{statement}
\begin{proof}
Every element outside the cyclic order-six kernel is
$r^as^b$ with $b$ odd and square $s^2\ne1$.
It has order four, whereas complex conjugation has order two.
For a nonconstant target the cocycle rule is instead
$c(g^2)=c(g)g(c(g))$, so that action must first be specified.
\end{proof}

\begin{statement}{D.14}{Proved}{The full lifting criterion and a class-field example}
For $K=\mathbb Q(\sqrt d)\ne\mathbb Q$, a surjective
constant $\operatorname{Dic}_3$ lift with quadratic field
$K$ exists exactly when $(d,-1)=0$.
Equivalently $K$ embeds in a cyclic quartic field.
An $S_3$-field with quadratic field $K$ always exists.
For $d=37$, take

$$
L=\operatorname{Spl}(X^3-X^2-3X+1),\qquad
 N=\text{the quartic subfield of }\mathbb Q(\zeta_{37}).
$$

Then $LN$ is such a dicyclic field, ramified only at
$2,37,\infty$, with $(e,f)=(3,4)$ at two and $(4,1)$
at $37$. Complex conjugation is the central involution.
\end{statement}
\begin{proof}
There is an isomorphism
$\operatorname{Dic}_3=S_3\times_{C_2}C_4$:
send $r$ to $((123),0)$ and $s$ to $((12),1)$.
An $S_3$-field and a cyclic quartic field with common $K$
intersect exactly in $K$, since the intersection degree
divides six and four. Their compositum gives the fibre product.
Class-field theory identifies the required data with an
order-four character of $\mathbb Q$ squaring to $\chi_K$
and a cubic character $\theta$ of $K$ satisfying
$\theta^\iota=\theta^{-1}$
\cite[V, Theorems 5.3,5.5]{D_MilneCFT}.
The first obstruction is the cyclic-four Bockstein
$\chi_K^2=(d,-1)$. A lift is automatically onto $C_4$.

For completeness the cubic character always exists.
Let $s=\dim_{\mathbb F_3}\operatorname{Cl}(K)/3$ and
choose $r=s+2$ rational primes splitting completely in
$K(\zeta_3)$, by Chebotarev
\cite[V, Theorem 3.23]{D_MilneCFT}.
For their product modulus $\mathfrak m$, put
$A=(\mathcal O_K/\mathfrak m)^\times$.
Each exchanged pair of residue factors contributes one
dimension to $\operatorname{Hom}(A,\mathbb F_3)^-$,
so that dimension is $r$.
The unit theorem gives
$\dim\mathcal O_K^\times/\mathcal O_K^{\times3}\le1$
\cite[Theorem 5.1]{D_MilneANT}.
If $A_0$ is the image in the ray class group $C_{\mathfrak m}$,
then $\dim\operatorname{Hom}(A_0,\mathbb F_3)^-\ge r-1$.
Applying $\operatorname{Hom}(-,\mathbb F_3)$ to
$0\to A_0\to C_{\mathfrak m}\to\operatorname{Cl}(K)\to0$
shows that the obstruction to extension lies in
$\operatorname{Ext}^1_{\mathbb Z}(\operatorname{Cl}(K),\mathbb F_3)$,
of dimension $s$.
Minus eigenspaces preserve exactness because two is invertible
modulo three. Thus at least $r-1-s=1$ anti-invariant characters
extend. A nonzero one gives a cyclic cubic extension of $K$
stable under conjugation, with inversion action.
A lift of conjugation has square fixed by inversion in $C_3$,
hence square one; the group is $S_3$.
This proves sufficiency of the sole quartic obstruction.

For $37$, the primitive quartic character is $\psi(2)=i$
modulo $37$, since two has order $36$.
Also $37=1^2+6^2$ exhibits the vanishing quaternion obstruction.
Let $\omega=(1+\sqrt{37})/2$.
Minkowski's bound is $\sqrt{37}/2<3.05$;
two is inert and the ideals of norm three are principal,
as $N((7+\sqrt{37})/2)=3$.
Thus the class number is one
\cite[Theorem 4.3]{D_MilneANT}.
The fundamental unit is $6+\sqrt{37}$, of norm minus one.
A smaller positive unit would have
$(a+b\sqrt{37})/2$ with $b=1$ or $2$ and
$a^2-37b^2=\pm4$; these equations give no smaller unit.
Division by its powers proves generation.
Both it and $-1$ reduce to one modulo two, so

$$
\operatorname{Cl}_{(2)}(K)=
 (\mathcal O_K/2)^\times=\mathbb F_4^\times.
$$

The character $\theta(\bar\omega)=\zeta_3$ is anti-invariant.

The cubic has no rational root and discriminant
$148=4\cdot37$, giving group $S_3$ and quadratic field $K$.
At two its translate is $X^3+2X^2-2X-2$, Eisenstein.
The local splitting field has tame inertia $C_3$ and
unramified quadratic subfield $K_2$.
At $37$ the reduction is $(X-4)^2(X-30)$,
and discriminant valuation one gives local splitting field
$K_{37}$. The positive discriminant gives three real roots.
Thus $L/K$ is precisely the cubic ray extension of modulus $(2)$.

For an explicit equation of $N$, let $H=\langle2^4\rangle$
modulo $37$ and $\eta_j=\sum_{h\in H}\zeta_{37}^{2^jh}$.
Exact multiplication yields

$$
\prod_{j=0}^3(T-\eta_j)=T^4+T^3+5T^2+7T+49.
$$

The periods are distinct: equality would give a nonzero
polynomial of degree at most $36$ with zero constant term,
divisible by $\Phi_{37}$, impossible since its coefficients
would all be equal.
At two the quartic Frobenius has order four, so adjoining it
to $L_2$ gives $(e,f)=(3,4)$.
At $37$ it is totally ramified quartic and contains $K_{37}$.
Finally $-1=2^{18}$ maps to two in $C_4$, while $L$ is real.
The discriminant, periods, root order and local reductions
are all included in the executed certificate.
\end{proof}

\begin{statement}{D.15}{Proved}{An explicit dihedral field}
For $\theta^4=-3$, the field $\mathbb Q(\theta,i)$ has group
$D_8$, quadratic subfields
$\mathbb Q(i),\mathbb Q(\sqrt{-3}),\mathbb Q(\sqrt3)$,
and only finite ramified primes two and three.
At both, $(e,f)=(4,2)$.
\end{statement}
\begin{proof}
At three the quartic is Eisenstein and tame, while adjoining
$i$ is unramified quadratic. The global degree is therefore
eight. The automorphisms
$a(\theta)=i\theta,a(i)=i$ and
$b(\theta)=\theta,b(i)=-i$ satisfy the dihedral relations.
The quadratic fields follow from $\theta^2$, $i$, and their
product. Discriminants exclude all other finite primes.

At two let $K_0=\mathbb Q_2(s)$, $s^2=-3$, the unramified
quadratic field, and $\omega=(1+s)/2$, so $\omega^2=\omega-1$.
The local field is $K_0(\sqrt{-1},\sqrt{s})$.
The unit squares modulo four are exactly
$1,-1+\omega,-\omega$, by enumerating the twelve unit residues.
None equals $-1=3$, $s=3+2\omega$, or $-s=1+2\omega$.
These three nontrivial square classes are independent as
needed for a biquadratic extension.
None is the unramified quadratic class: that class has
a unit representative $1+4a$, by lifting an irreducible
Artin--Schreier polynomial, and is square modulo four.
Thus all three quadratic subextensions ramify.
A nontrivial unramified quotient would yield an unramified
quadratic subextension, so the relative degree-four
extension is totally ramified. This proves $(4,2)$ at two.
\end{proof}

\begin{statement}{D.16}{Proved}{All central double covers of the exact M}
Let $N=M^{00}=\ker(e,t)$, $V=\mu_2^2$, and $\delta=\det f$.
Let $\operatorname{CExt}$ classify central extensions
with the identifications of kernel and quotient fixed.
Then

$$
\operatorname{CExt}_{\overline{\mathbb Q}}(M,\mu_2)=\mathbb F_2^5,
 \qquad
 \operatorname{CExt}_{\mathbb Q}(M,\mu_2)
   =\mathbb F_2^4\oplus(\mathbb Q^\times/\mathbb Q^{\times2})^2,
$$


$$
\boxed{\operatorname{CExt}_{\mathbb Z}(M,\mu_2)=\mathbb F_2^3.}
$$

The eight integral representatives are

$$
\widetilde M_{a,b,c}(R)=
 \{(m,z):z\in R^\times,\ z^2=\delta(m)^a e(m)^b t(m)^c\},
 \qquad a,b,c\in\{0,1\},
$$

with componentwise multiplication. They are finite locally
free of rank two over $M$ and fppf surjective, although
not necessarily surjective on $R$-points.
\end{statement}
\begin{proof}
Over $k=\overline{\mathbb Q}$ choose eigenvalues $z$ for $x$
and $u$ for $y$. On $N$, the other eigenvalues are
$\delta/z,\delta/u$, giving
$N_k=\operatorname{GL}_{2,k}\times\mathbb G_m^2$ as groups.
Its ring is a polynomial localization, so has trivial
Picard group, and its units modulo constants are freely
generated by $\delta,u,z$.
The fppf Kummer sequence
\cite[Lemma 59.28.3]{D_StacksKummer} gives
$H^1(N_k,\mu_2)=\langle\delta,u,z\rangle_{\mathbb F_2}$.
Each class is a square-root torsor of a character, with
its standard group law and chosen identity.
A pointed torsor isomorphism is unique, since a morphism
$N_k\to\mu_2$ is constant; it is multiplicative, since
its multiplication defect on $N_k^2$ is constant and
one at the identity.

The $e$-component sends $[u]$ to $[\delta]+[u]$ and fixes
$[\delta],[z]$; the $t$-component is trivial.
Thus restrictions of central extensions of $M_k$ lie
in $\langle\delta,z\rangle$. Both extend globally as
character square-root covers, since the first projection
always takes values in the commutative $T_\rho$.
After subtracting them an extension splits on $N_k$.
The section is unique because $\operatorname{Hom}(N_k,\mu_2)=0$;
conjugation therefore preserves its image.
Its right action gives descent of the finite flat torsor
along $M_k\to V_k$, including multiplication.
Thus the extension is the pullback of a unique extension
of $V_k$, and equivalences descend too.
D.9 gives the remaining $\mathbb F_2^3$.
This proves the geometric formula.

Over $\mathbb Q$, conjugation of $\mathbb Q(i)$ sends
$[u]\mapsto[\delta]+[u]$, while that of
$\mathbb Q(\sqrt{-3})$ sends $[z]\mapsto[\delta]+[z]$.
These distinct quadratic fields leave only $[\delta]$
invariant. Pointed automorphism groups of connected
extensions are trivial, so there are no additional
descent forms there. Subtracting the determinant cover
reduces to extensions of $V_{\mathbb Q}$.
Their eight constant geometric representatives have the
law in D.9, and automorphisms fixing kernel and quotient are

$$
(w,u,v)\longmapsto(w+\ell_1u+\ell_2v,u,v).
$$

Rational forms are therefore specified by two quadratic
Galois characters. The explicit formulas show that Baer
sum adds the three geometric parameters and these two
characters. Adding the determinant class proves the
rational formula.

For the integral argument, $N$ is smooth with geometrically
connected fibres. Indeed the norm $a^2+ab+b^2$ has derivatives
$2a+b,a+2b$ which cannot vanish together on its unit locus.
At three their common zero forces $a=b$ and norm zero.
The norm fibre product is therefore smooth.
Its geometric scalar factors are $\mathbb G_m^3$
away from two and three, and $\mathbb G_m^2\times\mathbb G_a$
at each exceptional characteristic: one algebra becomes
dual numbers, the other splits and permits eliminating
one coordinate. The factor $\operatorname{SL}_2$ preserves
connectedness. Hence $N$ is normal and integral.

Subtract the generic determinant class from an integral
extension. The unique splitting over $N_{\mathbb Q}$
extends over $N$: on a normal affine ring $B$, the finite
covering algebra maps under the generic section to
$\operatorname{Frac}B$, and each image is integral over
$B$, hence lies in $B$. Uniqueness glues the sections.
Multiplicativity and normality extend by flatness and
generic density on $N^2$ and $\widetilde M\times N$.
Affine descent along the faithfully flat quotient
$M\to\mu_2^2$ reduces to

$$
1\to\mu_2\to E\to\mu_2^2\to1.
$$


Such $E$ is commutative. Its commutator is a bilinear
pairing $V\times V\to\mu_2$, hence a morphism
$V\to\underline{\operatorname{Hom}}(V,\mu_2)=V^D
=(C_2^2)_{\mathbb Z}$.
The total scheme $V$ is connected: its four closed
sign sections cover its underlying space and meet at
the sole point above two. A morphism to this constant
scheme is constant, and the identity forces zero.
The commutator vanishes.

Exact Cartier duality
\cite[Sections 3.2.2--3.2.3 and Proposition 29]{D_Stix}
now gives

$$
0\to(C_2^2)_{\mathbb Z}\to E^D\to(C_2)_{\mathbb Z}\to0.
$$

The middle scheme is finite etale, since its map to the
quotient is a torsor under a finite etale group.
Every finite etale $\mathbb Z$-scheme is constant:
a connected component would be a number field unramified
at all finite primes, excluded except for $\mathbb Q$
by Minkowski \cite[Theorem 4.9]{D_MilneANT}.
Thus the problem is ordinary abelian extensions
$0\to C_2^2\to A\to C_2\to0$.
For a lift $r$ of the quotient generator, $2r\in C_2^2$
is independent of the lift and completely classifies it.
The four possibilities dualize to the covers $z^2=e^bt^c$.
The determinant class is independent already on $N_k$.
This proves the eight marked integral classes.

Finally $(a,b,c,d,e,t)=(1,2,1,1,1,1)$ with
$f=\operatorname{diag}(2,1)$ is a point of $N(\mathbb F_5)$.
It has no lift to $z^2=\delta$, since two is nonsquare.
This explicitly distinguishes fppf and point surjectivity.
\end{proof}

\begin{statement}{D.17}{Proved}{Generation and exact indices}
For every $n\ge1$, $\rho,\sigma$ generate
$\operatorname{SL}_2(\mathbb Z/n)$.
Adding the exchange matrix $\tau$ gives
$\det^{-1}(\{1,-1\})$ in GL, of index one for $n=1,2$
and $\varphi(n)/2$ otherwise.
The integral SL closure is all profinite SL.
The integral GL closure is $\det^{-1}(\{1,-1\})$,
of infinite index in profinite GL.
\end{statement}
\begin{proof}
The products $\rho\sigma$ and $\sigma\rho$ are the lower
and upper elementary matrices of parameters $1,-1$.
Their powers give every parameter modulo $n$.
Local row reduction and the elementary factorization of
$\operatorname{diag}(a,a^{-1})$ generate SL over
$\mathbb Z/p^r$. Chinese remaindering gives all $n$.
Since $\det\tau=-1$, adding it gives exactly the stated
determinants and indices. Finite quotients determine
the closures. The two-element determinant subgroup is
not open in $\widehat{\mathbb Z}^{\times}$.
The program checks every $1\le n\le48$.
An unrestricted ambient-group assertion is false:
embedding both generators in the first factor of
$\operatorname{SL}_2\times\mathbb G_m$ never generates
the second factor.
\end{proof}

\begin{statement}{D.18}{Proved}{A precise motivic arithmetic representation}
For an elliptic curve over a number field, torsion gives

$$
G_k\to\operatorname{Aut}(T_\ell E)
       =\operatorname{GL}_2(\mathbb Z_\ell),\qquad\det=\chi_\ell.
$$

This is the specialization of the universal elliptic local
system and rationally the dual realization of $h^1(E)$.
Matrices require a basis and change by conjugation with it.
\end{statement}
\begin{proof}
The finite etale groups $E[\ell^a]$ have compatible Galois
actions; their inverse limit is free of rank two.
The Weil pairing gives
$\bigwedge^2T_\ell E=\mathbb Z_\ell(1)$.
These constructions and specialization are those of
\cite[II, Section 6; V, Sections 7--8]{D_MilneEC}.
A fine modular curve of level at least three supplies
the universal local system after adjoining the level
structure; the moduli stack avoids this auxiliary step.
The correspondence

$$
[\Delta_E]-[\{0\}\times E]-[E\times\{0\}]
$$

is the idempotent projector to degree one: the other
two terms are orthogonal projectors to degrees zero and
two. Its realization is $H^1$, dual to the rational
Tate module. For $y^2=x^3-x$, the complete count at five
is eight, giving trace $-2$ and polynomial $T^2+2T+5$.
D.3 concerns additional norm coordinates and does not
contradict this GL representation.
\end{proof}

\begin{statement}{D.19}{Known}{Specified unipotent realizations}
Finite lower-central quotients of the punctured elliptic
fundamental Lie algebra have the Hain--Matsumoto universal
mixed elliptic realizations. A rational or tangential
basepoint gives a genuine arithmetic action; otherwise
the intrinsic action is outer.
For finitely generated discrete groups, rational unipotent
completion extended to $\mathbb Q_\ell$ agrees with
the corresponding unipotent completions of the profinite
and pro-$\ell$ completions.
\end{statement}
\begin{proof}[Exact reductions]
Theorem 6.9 and Corollary 6.10 of \cite{D_HMelliptic}
give the specified tangential punctured-elliptic fibre and
its extension in the defined category $\mathrm{MEM}_2$.
Brackets and lower-central quotients inherit those
realizations. Specialization at a pointed arithmetic
curve gives its fundamental-group section; forgetting
the point leaves the outer action.
The completion comparison is
\cite[Theorem A.6]{D_HMweighted}; the based application
is Remark 8.1 and Corollary A.11.
The crystalline theorem 9.8 there concerns the outer
action on the thrice-punctured line, with its stated
hypotheses. It does not apply to arbitrary integral
representations.
For a concrete quotient the free Lie algebra on $A,B$
has bases $A,B$; $[A,B]$;
$[A,[A,B]],[B,[A,B]]$ in degrees one to three,
of dimensions $2,1,2$.
Degree two is $\bigwedge^2H$, with cyclotomic action
in the covariant Tate convention.
\end{proof}

\begin{statement}{D.20}{Refuted}{An integral non-geometric representation}
The assignments

$$
r(\sigma)=\sigma,\qquad
 r(\rho)=P\rho P^{-1}=\begin{pmatrix}-1&3\\-1&2\end{pmatrix},
 \quad P=\begin{pmatrix}1&1\\0&1\end{pmatrix}
$$

define an integral modular representation which is not
of geometric origin on any finite-level modular curve.
This refutes an unrestricted geometric-origin reading
of C3-489, not every possible use of the word motivic.
\end{statement}
\begin{proof}
The squared and cubed generators both give $-I$, so the
amalgam defines the representation.
Its cusp generator has image

$$
T=\begin{pmatrix}1&-2\\-1&3\end{pmatrix},
 \qquad\det(X-T)=X^2-4X+1.
$$

The eigenvalues $2\pm\sqrt3$ are not roots of unity,
nor are any positive powers.
A level-$N$ cusp gives $T^N$, up to conjugacy and inversion.
The monodromy theorem
\cite[Theorem 1.2]{D_EsnaultKerz} requires
quasi-unipotent boundary monodromy for geometric-origin
local systems. This excludes the representation and
every finite pullback.
\end{proof}

\begin{statement}{D.21}{Proved}{Continuous arithmetic descent}
Let $\Pi=N\rtimes G_k$ be continuously split and
$r:N\to\operatorname{GL}(V)$ continuous, with $V$
finite-dimensional over a specified topological field $F$.
It extends continuously exactly when there is a continuous
homomorphism $A:G_k\to\operatorname{GL}(V)$ satisfying

$$
A_gr(n)A_g^{-1}=r(gng^{-1}).
$$

If $r$ is absolutely irreducible and continuous normalized
intertwiners $B_g$ exist, the scalar defect
$B_gB_hB_{gh}^{-1}$ defines a class in
$H^2_{\mathrm{cts}}(G_k,F^\times)$ with trivial scalar
action. Extension is equivalent to its vanishing.
Pointwise intertwiners and a continuous choice are prior
conditions.
\end{statement}
\begin{proof}
An extension restricts along the section to $A$.
Conversely $R(n,g)=r(n)A_g$ is a continuous representation
by direct multiplication.
For an irreducible $r$, Schur's lemma makes the defect
scalar. Associativity gives the cocycle identity,
and continuous scalar rescaling changes it by a
continuous coboundary. Thus vanishing is necessary
and sufficient under the stated existence hypothesis.

For an explicit earlier obstruction, let $N=C_2^2$
with the quotient involution swapping coordinates, and
$r(u,v)=(-1)^u$ over characteristic not two.
The swapped character differs at $(1,0)$, so no scalar
intertwiner exists. Conversely an invariant character
extends with $A_g=1$, and can be twisted by any continuous
character of $G_k$. D.18 provides the geometric positive
case with its intertwiners.
\end{proof}
